% id: 83-15
% book: 베이직쎈 대수 · 05 지수함수와 로그함수의 활용
% section: 유형 044 지수방정식의 실생활에의 활용
% figure: none
% answer: 
% answer_source: 
% variant_level: 0 (원본 전사)
% 이 파일은 build.mjs 가 items.json 에서 만든다. 고칠 때는 items.json 을 고친다.
\smash{\raisebox{1.5mm}{\dmkichul{베이직쎈 대수 83쪽 15번}}}
방사성 탄소 동위 원소 $^{14}\mathrm{C}$는 5730년마다 그 양이 반으로 줄어든다고 한다. 즉 처음 $^{14}\mathrm{C}$의 양이 $a\mkern3mu \mathrm{mg}$이었을 때 $x$년 후에 남아 있는 양을 $f(x)\mkern3mu \mathrm{mg}$이라 하면 다음과 같은 관계식이 성립한다고 한다.\par\smallskip\dispeq{$f(x)=a\cdot\left(\dfrac{1}{2}\right)^{\frac{x}{5730}}$}\par\smallskip\noindent 어떤 유물을 발굴하여 조사하였더니 $^{14}\mathrm{C}$가 $50\mkern3mu \mathrm{mg}$ 남아 있었다. 처음 $^{14}\mathrm{C}$의 양이 $400\mkern3mu \mathrm{mg}$이었다면 이 유물은 몇 년 전의 것이라 할 수 있는지 구하시오.
